% DFC/MANTIS v2.1 manuscript insertion.
% Evidence source: docs/VALIDATION_V21.md and immutable validation summaries.

\subsection{Full-FP32 Continual Adaptation on a Pretrained LLM}
\label{sec:qwen_partial}

We next cross the boundary from small continual-learning models to an actual
pretrained language model.  We use \texttt{Qwen/Qwen2.5-0.5B} at pinned Hub
revision \texttt{060db6499f32faf8b98477b0a26969ef7d8b9987} and sequentially
adapt the final transformer block with ordinary full-FP32 AdamW state.  The
trainable subspace contains 14,912,384 parameters.  External DER++ and
DFC-SIGN+DER++ receive exactly the same 512-byte external allocation, four
sequential tasks, 128 optimizer steps per task, identical two-prompt update
shape, and the same FP32 first- and second-moment allocations.  The only
additional addressability available to DFC-SIGN is the sign fiber already
present in the ordinary FP32 second moment: one exact bit per coordinate,
or 1,864,048 addressable bytes.  Replay records occupy 104 bytes.

The sealed evaluation uses five seeds (1901, 1931, 1951, 1987, and 2017).
All ten method--seed jobs completed and emitted immutable result rows.  The
aggregate-only workflow initially failed before analysis because its fresh
runner lacked NumPy; a recovery workflow downloaded the original artifacts,
recomputed every metric from the stored accuracy matrices, re-audited all
resource identities, and applied the original predeclared gate without
rerunning any learning job.  Table~\ref{tab:qwen_partial} reports the recovered
five-seed means.

\begin{table}[t]
\centering
\caption{Full-FP32 Qwen2.5-0.5B sequential partial fine-tuning under the
same 512-byte external envelope.  DFC differences are paired against external
DER++.  Accuracy differences are percentage points.}
\label{tab:qwen_partial}
\setlength{\tabcolsep}{4.5pt}
\begin{tabular}{lrrr}
\hline
Metric & External DER++ & DFC-SIGN+DER++ & $\Delta$ \\
\hline
Final average accuracy & 36.250\% & 79.375\% & +43.125 \\
Average forgetting & 84.167 & 25.833 & -58.333 \\
Current-task accuracy & 99.375\% & 98.125\% & -1.250 \\
Replay-record capacity & 4 & 17,928 & +17,924 \\
Mean wall time (s) & 252.31 & 242.90 & -9.41 \\
\hline
\end{tabular}
\end{table}

The paired final-accuracy gains are +34.375, +53.125, +50.000, +50.000, and
+28.125 points, so the effect is not carried by one seed.  The corresponding
forgetting differences are -45.833, -70.833, -70.833, -62.500, and -41.667
points.  Current-task differences are -3.125, 0, -3.125, +3.125, and -3.125
points.  The predeclared gate required mean final-accuracy gain of at least
+10 points, mean forgetting reduction of at least 5 points, and mean
current-task difference no worse than -5 points; all three conditions pass.
Thus the gain is a retention effect under a fixed external physical envelope,
not a collapse of current-task learning.

This experiment is deliberately described as a controlled sequential-task
partial-fine-tuning validation rather than a natural-domain benchmark.  A
separate five-seed AG News LoRA stress test provides a useful falsifier.  There,
DFC-SIGN+DER++ increases final average accuracy from 27.734\% to 37.109\% and
reduces forgetting from 82.500 to 61.667 points, but current-task accuracy
falls from 89.609\% to 79.219\%.  The predeclared plasticity condition is
therefore violated.  We retain this negative row rather than tune it away:
it shows that extra in-state replay capacity does not by itself remove the
stability--plasticity tradeoff for every adaptation protocol.

\subsection{Matched Fused Triton Deployment Path}
\label{sec:triton_v21}

The deployment artifact now contains two matched Triton primitives: a fused
reference AdamW kernel and a fused DFC-SIGN AdamW kernel.  They use the same
ordered FP32 AdamW arithmetic and the same parameter, gradient, first-moment,
and second-moment load/store footprint.  DFC-SIGN adds only sign extraction,
magnitude decoding, and sign re-embedding around that update.  The benchmark
first checks bitwise parameter and first-moment equality, decoded-second-moment
equality, and payload persistence; it then times both fused kernels with CUDA
Events for 1M, 4M, 16M, and 32M coordinates and reports latency, effective
bandwidth, and relative overhead.

At the time of this release, no CUDA-capable runner available to the artifact
has completed the hardware gate.  We therefore make no measured GPU-overhead
claim here.  The fused-vs.-fused benchmark, exactness guard, and workflow are
executable, and any future hardware number is admitted only from the produced
JSON artifact.  This preserves the evidence-tier rule that proposed or queued
hardware paths are code, not empirical rows.
